\begin{proof}[Proof of \cref{obj:lem:rough-full-record-sharp-lower}]
Fix an admissible tuple \(v=(p,\alpha,\beta,\gamma)\) with \(F(p)<1\), and an integer \(n\ge2\). Use the quantities of \cref{obj:def:sharp-frontier-scales}, in particular
\[
q=\frac{p-1}{p},
\qquad S=\alpha+\beta,
\qquad
D_p=1+2\alpha+\frac{\beta}{q}+\frac{S}{2\gamma}.
\]
Admissibility gives
\[
1<p\le2,\qquad \alpha,\beta>0,\qquad \gamma\ge\frac14,
\]
so
\[
0<q\le\frac12,\qquad S>0,\qquad
2\gamma+q>0,\qquad D_p>1.
\]
We first establish the original-model conclusions, splitting at \(N_{\mathrm{low}}(v)\).

\begin{enumerate}
\item \textbf{The small-sample branch.}
Suppose \(n<N_{\mathrm{low}}(v)\). Subtracting the two exponent formulas gives
\[
\begin{aligned}
E_4(p)-E_0(p)
&=
\frac{2S(2\gamma+q)-2\gamma qD_p}
     {(2\gamma+q)D_p}\\
&=
\frac{2\gamma q}{(2\gamma+q)D_p}
\left(
\frac{2S}{q}-2\alpha-\frac{\beta}{q}
+\frac{S}{2\gamma}-1
\right)\\
&=
\frac{2\gamma q}{(2\gamma+q)D_p}\{F(p)-1\}<0.
\end{aligned}
\]
The last identity uses \(q^{-1}-1=(p-1)^{-1}\). Thus \(E_0(p)-E_4(p)>0\). The definition of \(c_v\), together with \(n\le N_{\mathrm{low}}(v)\), yields
\[
\begin{aligned}
c_vn^{-E_4(p)}
&\le
c_0^{\mathrm e}
N_{\mathrm{low}}(v)^{-(E_0(p)-E_4(p))}
n^{-E_4(p)}\\
&\le
c_0^{\mathrm e}
n^{-(E_0(p)-E_4(p))}n^{-E_4(p)}
=
c_0^{\mathrm e}n^{-E_0(p)}.
\end{aligned}
\]

The selected priors are the paired-tent priors of
\cref{obj:def:converse-handle}, at the even rank
\[
N=2\left\lceil n^{2q/(2\gamma+q)}\right\rceil,
\qquad h=\frac1N,
\qquad \varepsilon_{n,v}=h^{\gamma/q},
\qquad L_{n,v}=h^{-\gamma/(p-1)}.
\]
We identify their canonical support laws explicitly. For
\(\sigma\in\{-1,1\}^{N/2}\), define the paired-tent field by
\[
g^{\mathrm{tent}}_\sigma(x)=
\begin{cases}
\sigma_j b_\triangle(Nx-(2j-2)),
&x\in[(2j-2)/N,(2j-1)/N],\\
-\sigma_j b_\triangle(Nx-(2j-1)),
&x\in[(2j-1)/N,2j/N],
\end{cases}
\qquad 1\le j\le N/2,
\]
where \(b_\triangle\) is the triangular tent of
\cref{obj:def:copula-frame-priors}; boundary values are zero.
Let \(P^{\mathrm{tent}}_\sigma\) be the law with uniform covariate,
propensity \(1/2\), and conditional outcome kernels
\[
Q_{0,P^{\mathrm{tent}}_\sigma}(\cdot\mid x)=\delta_0,
\]
\[
\begin{aligned}
Q_{1,P^{\mathrm{tent}}_\sigma}(\cdot\mid x)
={}&(1-\varepsilon_{n,v})\delta_0\\
&+\frac{\varepsilon_{n,v}}2
\left(1+\frac{g^{\mathrm{tent}}_\sigma(x)}{16}\right)\delta_{L_{n,v}}
+\frac{\varepsilon_{n,v}}2
\left(1-\frac{g^{\mathrm{tent}}_\sigma(x)}{16}\right)\delta_{-L_{n,v}}.
\end{aligned}
\]
Here \(\delta_y\) denotes a point mass at \(y\). Define \(P_0\) by the same kernels with the tent field replaced by zero. These are exactly the null and alternative laws selected by the paired-tent branch: its treated mean is
\[
\tau_{P^{\mathrm{tent}}_\sigma}(x)
=\frac{\varepsilon_{n,v}L_{n,v}}{16}g^{\mathrm{tent}}_\sigma(x)
=\frac{h^\gamma}{16}g^{\mathrm{tent}}_\sigma(x),
\qquad \tau_{P_0}(x)=0,
\]
since \(q^{-1}-(p-1)^{-1}=1\). Thus the selected alternative prior and its complete-record mixture are
\[
\pi_{1,n,v}
=2^{-N/2}\sum_{\sigma\in\{-1,1\}^{N/2}}
\delta_{P^{\mathrm{tent}}_\sigma},
\qquad
\overline P_1^{(n)}
=2^{-N/2}\sum_{\sigma\in\{-1,1\}^{N/2}}
(P^{\mathrm{tent}}_\sigma)^{\otimes n}.
\]
Before pushing forward to laws, each of the \(2^{N/2}\) sign configurations has weight \(2^{-N/2}\); hence
\[
2^{-N/2}>0,
\qquad
2^{N/2}2^{-N/2}=1.
\]
Every null sign configuration produces the same zero-effect law \(P_0\) just defined. Consequently,
\[
\pi_{0,n,v}=\delta_{P_0},
\qquad
\mathbb P_{0,n,v}=P_0^{\otimes n},
\qquad
\mathbb P_{1,n,v}=\overline P_1^{(n)}.
\]
Here the equality of priors is as probability measures on laws; the finite sign enumeration may retain repeated copies of \(P_0\).

We now certify these canonical laws directly. Ceiling bounds give
\[
2n^{2q/(2\gamma+q)}\le N\le4n^{2q/(2\gamma+q)},
\qquad N\ge2,
\qquad 0<h\le\frac12.
\]
Hence \(0<\varepsilon_{n,v}<1\) and \(L_{n,v}>0\). The displayed kernels have nonnegative masses summing to one, and their arm moments satisfy
\[
\int |y|^p\,Q_{0,P}(dy\mid x)=0,
\qquad
\int |y|^p\,Q_{1,P}(dy\mid x)
=\varepsilon_{n,v}L_{n,v}^p=1
\]
for \(P=P_0\) and every \(P=P^{\mathrm{tent}}_\sigma\).
The propensity and baseline mean are constant, while
\[
\|g^{\mathrm{tent}}_\sigma\|_\infty\le1,
\qquad
|g^{\mathrm{tent}}_\sigma(x)-g^{\mathrm{tent}}_\sigma(z)|
\le8N|x-z|.
\]
For the latter bound, each paired bump has Lipschitz constant at most \(4N\), and at most two paired bumps contribute at the two endpoints. Combining this bound with the envelope bound gives, for \(0<\gamma\le1\),
\[
h^\gamma
|g^{\mathrm{tent}}_\sigma(x)-g^{\mathrm{tent}}_\sigma(z)|
\le8|x-z|^\gamma.
\]
Indeed, when \(N|x-z|\le1\), use
\(N|x-z|\le(N|x-z|)^\gamma\); when \(N|x-z|>1\), use the envelope bound \(2\le8(N|x-z|)^\gamma\).
Consequently,
\[
\|\tau_{P^{\mathrm{tent}}_\sigma}\|_\infty\le\frac1{16},
\qquad
|\tau_{P^{\mathrm{tent}}_\sigma}(x)
-\tau_{P^{\mathrm{tent}}_\sigma}(z)|
\le\frac12|x-z|^\gamma.
\]
These calculations verify all restrictions of
\cref{obj:def:model,obj:def:null}, so
\[
P_0\in H_0(v),
\qquad P^{\mathrm{tent}}_\sigma\in\mathcal M_v.
\]
Opposite pairing and integration of the squared triangular tent give
\[
\int g^{\mathrm{tent}}_\sigma\,d\lambda=0,
\qquad
\int (g^{\mathrm{tent}}_\sigma)^2\,d\lambda=\frac13.
\]
Thus the distance of each canonical alternative is
\[
d(P^{\mathrm{tent}}_\sigma)
=\frac{h^\gamma}{16\sqrt3}
\ge\frac{4^{-\gamma}}{16\sqrt3}
 n^{-2\gamma q/(2\gamma+q)}
=c_0^{\mathrm e}n^{-E_0(p)},
\qquad
d(P^{\mathrm{tent}}_\sigma)\le\frac1{16}<d_0.
\]

To compare the canonical complete-record mixtures, define the one-record likelihood ratio on the full record space by
\[
\ell^{\mathrm{tent}}_\sigma(X,A,Y)
=1+\frac{g^{\mathrm{tent}}_\sigma(X)}{16}
\begin{cases}
1,&A=1,\ Y=L_{n,v},\\
-1,&A=1,\ Y=-L_{n,v},\\
0,&\text{otherwise}.
\end{cases}
\]
The conditional masses show directly that
\[
\frac{dP^{\mathrm{tent}}_\sigma}{dP_0}
=\ell^{\mathrm{tent}}_\sigma
\quad P_0\text{-almost surely}.
\]
For two sign configurations \(\sigma,\sigma'\), integration over all record categories and then over the paired cells gives
\[
\int g^{\mathrm{tent}}_\sigma
 g^{\mathrm{tent}}_{\sigma'}\,d\lambda
=\frac{2}{3N}\sum_{j=1}^{N/2}\sigma_j\sigma'_j,
\]
\[
\int\ell^{\mathrm{tent}}_\sigma
\ell^{\mathrm{tent}}_{\sigma'}\,dP_0
=1+\frac{\varepsilon_{n,v}}{768N}
\sum_{j=1}^{N/2}\sigma_j\sigma'_j.
\]
Write \(\chi^2(Q\Vert P)=\int(dQ/dP-1)^2\,dP\) for the directed chi-square divergence. Independence of the records and finite-mixture expansion yield
\[
1+\chi^2(\overline P_1^{(n)}\Vert P_0^{\otimes n})
=\mathbb E_{\sigma,\sigma'}
\left[
\left(1+\frac{\varepsilon_{n,v}}{768N}
\sum_{j=1}^{N/2}\sigma_j\sigma'_j\right)^n
\right],
\]
where both sign configurations are independent and uniform.
The expression inside the power is positive. Using
\((1+\vartheta_{\mathrm{exp}})^n\le e^{n\vartheta_{\mathrm{exp}}}\) for \(\vartheta_{\mathrm{exp}}>-1\), independence of the sign products, and
\(\cosh(\vartheta_{\mathrm{exp}})\le e^{\vartheta_{\mathrm{exp}}^2/2}\) for \(\vartheta_{\mathrm{exp}}\in\mathbb R\), we obtain
\[
\begin{aligned}
1+\chi^2(\overline P_1^{(n)}\Vert P_0^{\otimes n})
&\le
\left[\cosh\!\left(\frac{n\varepsilon_{n,v}}{768N}\right)\right]^{N/2}\\
&\le
\exp\!\left(\frac{n^2\varepsilon_{n,v}^2h}{36\cdot16^4}\right).
\end{aligned}
\]
The selected rank calibrates this exponent because
\[
n^2\varepsilon_{n,v}^2h
=n^2N^{-(1+2\gamma/q)}\le1.
\]
Here the lower rank bound implies
\(N\ge n^{2q/(2\gamma+q)}\), and
\(\frac{2q}{2\gamma+q}(1+2\gamma/q)=2\).
Since \(e^{\vartheta_{\mathrm{exp}}}\le(1-\vartheta_{\mathrm{exp}})^{-1}\) for \(0\le\vartheta_{\mathrm{exp}}<1\),
\[
1+\chi^2(\overline P_1^{(n)}\Vert P_0^{\otimes n})
\le\exp\!\left(\frac1{36\cdot16^4}\right)<2.
\]
Cauchy--Schwarz therefore gives
\[
\operatorname{TV}(\mathbb P_{0,n,v},\mathbb P_{1,n,v})
\le\frac12
\sqrt{\chi^2(\overline P_1^{(n)}\Vert P_0^{\otimes n})}
<\frac12.
\]
The model-supremum bound is supplied by
\cref{obj:lem:paired-tent-full-record-lower}:
\[
d_0\le D_v.
\]
Together with the multiplier comparison, these canonical calculations give, for every positive-weight alternative support law,
\[
P\in\mathcal M_v,
\qquad
c_vn^{-E_4(p)}
\le c_0^{\mathrm e}n^{-E_0(p)}
\le d(P)<d_0.
\]
Finally, for every size-valid test \(\phi\), averaging over its independent uniform seed and using the total-variation expectation bound gives
\[
\begin{aligned}
\sup_{\substack{P\in\mathcal M_v\\
d(P)\ge c_0^{\mathrm e}n^{-E_0(p)}}}
\{1-\mathsf R_n(P,\phi)\}
&\ge\int\{1-\mathsf R_n(P,\phi)\}\,\pi_{1,n,v}(dP)\\
&\ge\frac9{10}
-\operatorname{TV}(P_0^{\otimes n},\overline P_1^{(n)})
>\frac25.
\end{aligned}
\]
Taking the infimum over size-valid tests proves the canonical risk bound
\[
B_n\!\left(v,c_0^{\mathrm e}n^{-E_0(p)}\right)\ge\frac25.
\]
% lean: CausalSmith.Stat.FinitepHomogeneityDensegamma.tent_canonical_lower_receipt

For each test satisfying the size constraint, decreasing the separation enlarges the alternative set, so
\[
\sup_{\substack{P\in\mathcal M_v\\
d(P)\ge c_0^{\mathrm e}n^{-E_0(p)}}}
\{1-\mathsf R_n(P,\phi)\}
\le
\sup_{\substack{P\in\mathcal M_v\\
d(P)\ge c_vn^{-E_4(p)}}}
\{1-\mathsf R_n(P,\phi)\}.
\]
The positive-weight alternative supplied by the paired-tent construction makes both sets nonempty. Taking infima over the same size-valid tests and using its risk bound gives
\[
\frac25
\le B_n\!\left(v,c_0^{\mathrm e}n^{-E_0(p)}\right)
\le B_n\!\left(v,c_vn^{-E_4(p)}\right).
\]
Finally, \(n<N_{\mathrm{low}}(v)\) excludes the hypothesis
\(n\ge N_{\mathrm{low}}(v)\) of the additional rank and activity conclusions.
% lean: CausalSmith.Stat.FinitepHomogeneityDensegamma.rough_small_sample_lower

\item \textbf{Ranks in the large-sample branch.}
Suppose now \(n\ge N_{\mathrm{low}}(v)\). Define the selected fine and coarse ranks by
\[
K=
2^{\left\lceil\log_2(H_{\mathrm{fine}}n^{2/D_p})\right\rceil},
\qquad
M=
2^{\left\lfloor\log_2(K^{S/\gamma}/16)\right\rfloor}.
\]
Since \(p-1\le1\) and \(q\le1/2\),
\[
2S+\frac{S}{2\gamma}\le F(p)<1,
\qquad
S<\frac{2\gamma}{4\gamma+1}\le\gamma.
\]
Also,
\[
D_p-1
=
2\alpha+\frac{\beta}{q}+\frac{S}{2\gamma}
\le F(p)<1,
\qquad
1<D_p<2.
\]
In particular \(2/D_p\ge1\), and therefore \(n^{2/D_p}\ge n\).
The ceiling inequalities
\[
\log_2(H_{\mathrm{fine}}n^{2/D_p})
\le
\left\lceil\log_2(H_{\mathrm{fine}}n^{2/D_p})\right\rceil
\le
\log_2(H_{\mathrm{fine}}n^{2/D_p})+1
\]
give
\[
H_{\mathrm{fine}}n^{2/D_p}\le K
\le2H_{\mathrm{fine}}n^{2/D_p},
\qquad
2^{40}n\le K.
\]

The definition of the threshold gives
\[
n\ge32^{D_p\gamma/(2S)},
\qquad
K^{S/\gamma}\ge n^{2S/(D_p\gamma)}\ge32.
\]
Thus the coarse target \(K^{S/\gamma}/16\) is at least \(2\).
Applying the floor inequalities to its base-two logarithm gives
\[
\frac{K^{S/\gamma}}{32}\le M
\le\frac{K^{S/\gamma}}{16},
\qquad M\ge2.
\]
Both ranks are powers of two. Since \(K\ge1\) and \(S/\gamma<1\),
\[
16M\le K^{S/\gamma}\le K,
\qquad
M\le K^{S/\gamma}.
\]
These are all the required rank-legality conditions.
% lean: CausalSmith.Stat.FinitepHomogeneityDensegamma.rough_rank_legality

\item \textbf{Normalization, support, and separation.}
At these ranks, define the legality tuning and its copula priors by
\[
a_{\mathrm{leg}}=\frac{K^{-\alpha}}{16},
\qquad
b_{\mathrm{leg}}=\frac{K^{-\beta}}{256},
\qquad
u=\frac1{16},
\qquad
\varepsilon=(16b_{\mathrm{leg}})^{1/q},
\qquad
L_{\mathrm{mark}}=\varepsilon^{-1/p},
\]
\[
\pi_\nu=\pi^{K,M}_{\nu;v,a_{\mathrm{leg}},u,\varepsilon,L_{\mathrm{mark}}},
\qquad
\mathbb P_\nu=\int P^{\otimes n}\,\pi_\nu(dP),
\qquad \nu\in\{0,1\}.
\]
Because \(K\ge1\) and \(\alpha,\beta>0\),
\[
0<a_{\mathrm{leg}}\le\frac1{16},
\qquad
0<b_{\mathrm{leg}}\le\frac1{256},
\qquad
0<u\le\frac1{16},
\qquad
0<\varepsilon<1,
\qquad L_{\mathrm{mark}}>0.
\]

For completeness, the coefficient weights in
\cref{obj:def:copula-frame-priors} are normalized as follows.
The triangular tent satisfies
\[
0\le b_\triangle(t_{\triangle})=2\min\{t_{\triangle},1-t_{\triangle}\}\le1
\qquad(t_{\triangle}\in[0,1]),
\]
so \(|g_\sigma(x)|\le1\) on \([0,1]\). With \(\kappa=1/16\), each node-pair probability satisfies
\[
\frac{1+\nu\kappa g_\sigma(i/K)\lambda_i\eta_i}{4}
\ge\frac{1-\kappa}{4}>0,
\]
and cancellation of the sign product gives
\[
\sum_{(\lambda_i,\eta_i)\in\{-1,1\}^2}
\frac{1+\nu\kappa g_\sigma(i/K)\lambda_i\eta_i}{4}
=1.
\]
Conditional independence therefore makes the sum of the fine-pair product weights equal to one for every coarse-sign configuration. Summing over the fair coarse signs then gives
\[
\sum_{\sigma\in\{-1,1\}^{M/2}}2^{-M/2}=1.
\]
Thus both finite coefficient distributions, and their pushforwards \(\pi_0,\pi_1\), have nonnegative weights summing to one.
% lean: CausalSmith.Stat.FinitepHomogeneityDensegamma.copula_prior_normalized

The rank conditions established above allow us to apply
\cref{obj:lem:copula-frame-legality}. Every null support law belongs to \(H_0(v)\), and every alternative support law satisfies
\[
P\in\mathcal M_v,
\qquad
2^{-17}K^{-S}\le d(P)<d_0.
\]
The upper rounding bound and \(S>0\) imply
\[
\begin{aligned}
c_vn^{-E_4(p)}
&\le k_vn^{-E_4(p)}\\
&=2^{-17}(2H_{\mathrm{fine}})^{-S}n^{-2S/D_p}\\
&=2^{-17}(2H_{\mathrm{fine}}n^{2/D_p})^{-S}\\
&\le2^{-17}K^{-S}\le d(P).
\end{aligned}
\]
This proves the required alternative placement for the copula pair.
% lean: CausalSmith.Stat.FinitepHomogeneityDensegamma.rough_rank_separation

\item \textbf{Activity calibration.}
Define the deterministic first- and second-component activity bounds by
\[
\mathcal A_{\mathrm{det}}=
\frac{2^{38}a_{\mathrm{leg}}^4u^4\varepsilon^2n^2}{K},
\qquad
\mathcal B_{\mathrm{det}}=
\frac{2^{56}a_{\mathrm{leg}}^4u^4\varepsilon^2n^4}{MK^2}.
\]
In \cref{obj:lem:full-record-copula-component-bound}, let \(\mu\) be the common augmentation law and write \(\mathcal A_{\mathrm{cond}},\mathcal B_{\mathrm{cond}}\) for its conditional chi-square exponent budgets. The expectation bounds supplied there are
\[
\int\mathcal A_{\mathrm{cond}}\,d\mu
\le\mathcal A_{\mathrm{det}},
\qquad
\int\mathcal B_{\mathrm{cond}}\,d\mu
\le\mathcal B_{\mathrm{det}}.
\]
We calibrate the deterministic quantities on the right.

First, raising the lower fine-rank bound to \(D_p>0\) gives
\[
K^{D_p}\ge H_{\mathrm{fine}}^{D_p}n^2
\ge H_{\mathrm{fine}}n^2.
\]
Consequently,
\[
n^2K^{-D_p}\le\frac1{H_{\mathrm{fine}}},
\qquad
n^4K^{-2D_p}
=\bigl(n^2K^{-D_p}\bigr)^2
\le\frac1{H_{\mathrm{fine}}^2}.
\]

The tuning gives the exact amplitude product
\[
\varepsilon=16^{-1/q}K^{-\beta/q},
\qquad
a_{\mathrm{leg}}^4u^4\varepsilon^2
=
16^{-8-2/q}K^{-(4\alpha+2\beta/q)}
\le
2^{-48}K^{-(4\alpha+2\beta/q)},
\]
where the last inequality follows from \(q\le1/2\).
Furthermore,
\[
\frac{S}{2\gamma}\le2S
=2\alpha+2\beta
\le2\alpha+\frac{\beta}{q},
\]
and hence
\[
D_p\le1+4\alpha+\frac{2\beta}{q}.
\]
As \(K\ge1\), these inequalities yield
\[
\begin{aligned}
\mathcal A_{\mathrm{det}}
&\le
2^{-10}n^2K^{-(1+4\alpha+2\beta/q)}\\
&\le2^{-10}n^2K^{-D_p}
\le\frac{2^{-10}}{H_{\mathrm{fine}}}.
\end{aligned}
\]
For the second budget, use the lower coarse-rank bound:
\[
\begin{aligned}
\mathcal B_{\mathrm{det}}
&\le
\frac{2^8n^4K^{-(2+4\alpha+2\beta/q)}}{M}\\
&\le
2^{13}n^4
K^{-(2+4\alpha+2\beta/q+S/\gamma)}\\
&=
2^{13}n^4K^{-2D_p}
\le\frac{2^{13}}{H_{\mathrm{fine}}^2}.
\end{aligned}
\]
The exponent identity in the last line is
\[
2D_p=2+4\alpha+\frac{2\beta}{q}+\frac{S}{\gamma}.
\]
Since \(H_{\mathrm{fine}}=2^{40}\),
\[
\mathcal A_{\mathrm{det}}+\mathcal B_{\mathrm{det}}
\le
\frac{2^{-10}}{H_{\mathrm{fine}}}
+\frac{2^{13}}{H_{\mathrm{fine}}^2}
=
2^{-50}+2^{-67}<2^{-16}.
\]
All sample-size, rank, amplitude, and mark hypotheses of
\cref{obj:lem:full-record-copula-component-bound} have now been discharged. Its complete-record comparison gives
\[
\operatorname{TV}(\mathbb P_0,\mathbb P_1)
<\frac18<\frac12.
\]
% lean: CausalSmith.Stat.FinitepHomogeneityDensegamma.tail_rough_activity_budgets

\item \textbf{Testing risk and identification of the selected pair.}
Let \(\phi\) be any randomized test satisfying the size constraint in
\cref{obj:def:testing-risk}. Average over the independent public uniform seed \(U\) of
\cref{obj:def:original-record-experiment}, defining
\[
\overline\phi:
\bigl([0,1]\times\{0,1\}\times\mathbb R\bigr)^n
\longrightarrow[0,1],
\qquad
\overline\phi(\mathcal D_n)
=
\mathbb E_U[\phi(\mathcal D_n,U)].
\]
Finite-mixture linearity and null support give
\[
\int\overline\phi\,d\mathbb P_0
=
\int\mathsf R_n(P,\phi)\,\pi_0(dP)
\le\frac1{10}.
\]
For a measurable function taking values in \([0,1]\), the expectation difference is bounded by total variation. Thus
\[
\left|
\int\overline\phi\,d\mathbb P_1
-
\int\overline\phi\,d\mathbb P_0
\right|
\le\operatorname{TV}(\mathbb P_0,\mathbb P_1).
\]
Using alternative support and normalization,
\[
\begin{aligned}
\sup_{\substack{P\in\mathcal M_v\\
d(P)\ge c_vn^{-E_4(p)}}}
\{1-\mathsf R_n(P,\phi)\}
&\ge
\int\{1-\mathsf R_n(P,\phi)\}\,\pi_1(dP)\\
&=
1-\int\overline\phi\,d\mathbb P_1\\
&\ge
\frac9{10}-\operatorname{TV}(\mathbb P_0,\mathbb P_1).
\end{aligned}
\]
Normalization ensures that at least one alternative weight is positive, so the alternative set is nonempty; zero-weight terms contribute zero. The zero test satisfies the size constraint, so the infimum is also over a nonempty set. Taking that infimum gives
\[
B_n\!\left(v,c_vn^{-E_4(p)}\right)
\ge
\frac9{10}-\operatorname{TV}(\mathbb P_0,\mathbb P_1)
>\frac25.
\]
% lean: CausalSmith.Stat.FinitepHomogeneityDensegamma.finite_two_prior_testing_lower

It remains to identify this copula pair with the selected pair. In the present branch, the formulas in
\cref{obj:def:converse-handle} satisfy
\[
\varepsilon_{n,v}
=
16^{-1/q}K^{-\beta/q}
=
\left(16\frac{K^{-\beta}}{256}\right)^{1/q}
=\varepsilon,
\qquad
L_{n,v}=\varepsilon_{n,v}^{-1/p}=L_{\mathrm{mark}}.
\]
Therefore
\[
\pi_{\nu,n,v}=\pi_\nu,
\qquad
\mathbb P_{\nu,n,v}=\mathbb P_\nu,
\qquad \nu\in\{0,1\}.
\]
For the selected pair, the conditional budgets supplied by
\cref{obj:lem:full-record-copula-component-bound} consequently satisfy
\[
\int\mathcal A_{\mathrm{cond}}\,d\mu
\le\mathcal A_{\mathrm{det}}
\le\frac{2^{-10}}{H_{\mathrm{fine}}},
\qquad
\int\mathcal B_{\mathrm{cond}}\,d\mu
\le\mathcal B_{\mathrm{det}}
\le\frac{2^{13}}{H_{\mathrm{fine}}^2}.
\]
The calibrated activity bounds apply to the deterministic right-hand sides of the component expectation estimates.
Finally, the model witness supplied by
\cref{obj:lem:paired-tent-full-record-lower} gives
\[
d_0\le D_v.
\]
Together with the preceding calculations, this establishes normalization, support, separation, the model-supremum bound, testing risk, total variation, rank legality, the fine-rank bound, and the deterministic activity calibration in the large-sample branch.
% lean: CausalSmith.Stat.FinitepHomogeneityDensegamma.tail_rough_large_sample_lower

\item \textbf{The bounded small-sample branch.}
Fix an admissible \(w=(\alpha,\beta,\gamma)\), put
\[
v=(2,\alpha,\beta,\gamma),
\qquad
c_w^{\mathrm b}=c_v,
\qquad
E_0(2)=\frac{2\gamma}{4\gamma+1},
\]
and suppose \(F(2)<1\). This \(v\) is admissible. If \(n<N_{\mathrm{low}}(v)\), the exponent and multiplier calculation at \(p=2\) gives
\[
c_w^{\mathrm b}n^{-E_4(2)}
\le c_0^{\mathrm e}n^{-E_0(2)}.
\]
The bounded fallback in \cref{obj:synth_1} selects the signed-binary paired-tent package of
\cref{obj:lem:paired-tent-full-record-lower}:
\[
\pi^{\mathrm b}_{0,n,v}=\delta_{P_0^{\mathrm b}},
\qquad
\pi^{\mathrm b}_{1,n,v}=\pi_1^{\mathrm b},
\]
\[
\mathbb P^{\mathrm b}_{0,n,v}
=(P_0^{\mathrm b})^{\otimes n},
\qquad
\mathbb P^{\mathrm b}_{1,n,v}
=\overline P_{1,\mathrm b}^{(n)}.
\]
The point mass and the finite alternative prior are normalized. Signed-binary support implies \(|Y|\le1\), so their support laws satisfy
\[
P_0^{\mathrm b}\in H_0^{\mathrm b}(w),
\]
and, on every positive-weight alternative,
\[
P\in\mathcal M_w^{\mathrm b},
\qquad
c_w^{\mathrm b}n^{-E_4(2)}
\le c_0^{\mathrm e}n^{-E_0(2)}
\le d(P)<d_0.
\]
The same paired-tent bound supplies
\[
\operatorname{TV}
\bigl(\mathbb P^{\mathrm b}_{0,n,v},
      \mathbb P^{\mathrm b}_{1,n,v}\bigr)<\frac12,
\qquad
B_n^{\mathrm b}\!\left(w,c_0^{\mathrm e}n^{-E_0(2)}\right)
\ge\frac25.
\]
For every bounded-null size-valid test, the alternative-set inclusion gives
\[
\sup_{\substack{P\in\mathcal M_w^{\mathrm b}\\
d(P)\ge c_0^{\mathrm e}n^{-E_0(2)}}}
\{1-\mathsf R_n(P,\phi)\}
\le
\sup_{\substack{P\in\mathcal M_w^{\mathrm b}\\
d(P)\ge c_w^{\mathrm b}n^{-E_4(2)}}}
\{1-\mathsf R_n(P,\phi)\}.
\]
The paired-tent prior has a positive weight, ensuring nonempty alternative sets. Taking infima yields
\[
B_n^{\mathrm b}\!\left(w,c_w^{\mathrm b}n^{-E_4(2)}\right)
\ge
B_n^{\mathrm b}\!\left(w,c_0^{\mathrm e}n^{-E_0(2)}\right)
\ge\frac25.
\]
% lean: CausalSmith.Stat.FinitepHomogeneityDensegamma.bounded_rough_small_sample_lower

\item \textbf{The bounded large-sample branch.}
Suppose \(n\ge N_{\mathrm{low}}(v)\). Evaluate the rank formulas and amplitudes \(a_{\mathrm{leg}},b_{\mathrm{leg}}\) at \(v=(2,\alpha,\beta,\gamma)\). The rank calculation gives
\[
K=
2^{\left\lceil\log_2(H_{\mathrm{fine}}n^{2/D_2})\right\rceil},
\qquad
M=
2^{\left\lfloor\log_2(K^{S/\gamma}/16)\right\rfloor},
\]
\[
H_{\mathrm{fine}}n^{2/D_2}\le K
\le2H_{\mathrm{fine}}n^{2/D_2},
\qquad
2^{40}n\le K,
\]
\[
\frac{K^{S/\gamma}}{32}\le M
\le\frac{K^{S/\gamma}}{16},
\qquad
M\ge2,\qquad16M\le K,\qquad M\le K^{S/\gamma}.
\]
Use the separate bounded tuning
\[
a_{\mathrm{leg}}=\frac{K^{-\alpha}}{16},
\qquad
b_{\mathrm{leg}}=\frac{K^{-\beta}}{256},
\qquad
u_{\mathrm b}=2b_{\mathrm{leg}}=\frac{K^{-\beta}}{128},
\qquad
\varepsilon_{\mathrm b}=\frac12,
\qquad
L_{\mathrm b}=1,
\]
and define its priors and complete-record mixtures by
\[
\pi_\nu^{\mathrm{b,frame}}
=
\pi^{K,M}_{\nu;v,a_{\mathrm{leg}},u_{\mathrm b},\varepsilon_{\mathrm b},L_{\mathrm b}},
\qquad
\mathbb P_\nu^{\mathrm{b,frame}}
=
\int P^{\otimes n}\,\pi_\nu^{\mathrm{b,frame}}(dP),
\qquad \nu\in\{0,1\}.
\]
Here
\[
0<a_{\mathrm{leg}}\le\frac1{16},
\qquad
0<u_{\mathrm b}\le\frac1{128}\le\frac1{16},
\qquad
0<\varepsilon_{\mathrm b}<1,
\qquad L_{\mathrm b}>0.
\]
The coefficient weights depend on the ranks and the copula coupling, so the product-weight calculation proves normalization for both bounded priors as well. The bounded clause of
\cref{obj:lem:copula-frame-legality} supplies
\[
P\in H_0^{\mathrm b}(w)
\quad\text{on null support},
\]
\[
P\in\mathcal M_w^{\mathrm b},
\qquad
2^{-17}K^{-S}\le d(P)<d_0
\quad\text{on alternative support}.
\]
The upper fine-rank bound gives
\[
c_w^{\mathrm b}n^{-E_4(2)}
\le k_vn^{-E_4(2)}
=
2^{-17}(2H_{\mathrm{fine}}n^{2/D_2})^{-S}
\le2^{-17}K^{-S}.
\]

For this tuning, the exact amplitude product is
\[
a_{\mathrm{leg}}^4u_{\mathrm b}^4\varepsilon_{\mathrm b}^2
=2^{-46}K^{-4S}.
\]
Thus the deterministic component budgets, denoted by
\(\mathcal A_{\mathrm b}\) and \(\mathcal B_{\mathrm b}\), are
\[
\mathcal A_{\mathrm b}
=2^{-8}n^2K^{-(1+4S)},
\qquad
\mathcal B_{\mathrm b}
=\frac{2^{10}n^4K^{-(2+4S)}}{M}.
\]
At \(p=2\),
\[
D_2=1+2S+\frac{S}{2\gamma}
\le1+4S,
\]
because \(\gamma\ge1/4\). Raising the fine-rank bound to \(D_2\), and then squaring the resulting reciprocal bound, gives
\[
n^2K^{-D_2}\le\frac1{H_{\mathrm{fine}}},
\qquad
n^4K^{-2D_2}\le\frac1{H_{\mathrm{fine}}^2}.
\]
It follows that
\[
\mathcal A_{\mathrm b}
\le2^{-8}n^2K^{-D_2}
\le\frac{2^{-8}}{H_{\mathrm{fine}}},
\]
and, using \(M\ge K^{S/\gamma}/32\),
\[
\begin{aligned}
\mathcal B_{\mathrm b}
&\le
2^{15}n^4K^{-(2+4S+S/\gamma)}\\
&=
2^{15}n^4K^{-2D_2}
\le\frac{2^{15}}{H_{\mathrm{fine}}^2}.
\end{aligned}
\]
Consequently,
\[
\mathcal A_{\mathrm b}+\mathcal B_{\mathrm b}
\le
\frac{2^{-8}}{H_{\mathrm{fine}}}
+\frac{2^{15}}{H_{\mathrm{fine}}^2}
=
2^{-48}+2^{-65}<2^{-16}.
\]
Applying \cref{obj:lem:full-record-copula-component-bound}, with the verified bounded parameters, gives
\[
\operatorname{TV}
\bigl(\mathbb P_0^{\mathrm{b,frame}},
      \mathbb P_1^{\mathrm{b,frame}}\bigr)
<\frac18<\frac12.
\]
% lean: CausalSmith.Stat.FinitepHomogeneityDensegamma.bounded_rough_activity_budgets

For a test satisfying the bounded-null size constraint, seed averaging, null support, and the total-variation expectation bound give
\[
\int\overline\phi\,d\mathbb P_0^{\mathrm{b,frame}}\le\frac1{10},
\]
\[
\begin{aligned}
\sup_{\substack{P\in\mathcal M_w^{\mathrm b}\\
d(P)\ge c_w^{\mathrm b}n^{-E_4(2)}}}
\{1-\mathsf R_n(P,\phi)\}
&\ge
1-\int\overline\phi\,d\mathbb P_1^{\mathrm{b,frame}}\\
&\ge
\frac9{10}
-\operatorname{TV}
\bigl(\mathbb P_0^{\mathrm{b,frame}},
      \mathbb P_1^{\mathrm{b,frame}}\bigr).
\end{aligned}
\]
Both priors are normalized, so their alternative support is nonempty; the zero test again satisfies the size constraint. Taking the infimum therefore yields
\[
B_n^{\mathrm b}\!\left(w,c_w^{\mathrm b}n^{-E_4(2)}\right)
\ge
\frac9{10}
-\operatorname{TV}
\bigl(\mathbb P_0^{\mathrm{b,frame}},
      \mathbb P_1^{\mathrm{b,frame}}\bigr)
>\frac25.
\]
Finally, the selection in \cref{obj:synth_1} identifies
\[
\pi^{\mathrm b}_{\nu,n,v}=\pi_\nu^{\mathrm{b,frame}},
\qquad
\mathbb P^{\mathrm b}_{\nu,n,v}
=\mathbb P_\nu^{\mathrm{b,frame}},
\qquad \nu\in\{0,1\}.
\]
This transfers normalization, support, separation, total variation, and the testing-risk bound to the selected bounded pair, completing the bounded clause.
% lean: CausalSmith.Stat.FinitepHomogeneityDensegamma.bounded_rough_large_sample_lower
\end{enumerate}
% lean: CausalSmith.Stat.FinitepHomogeneityDensegamma.rough_full_record_sharp_lower
\end{proof}
